% =====================================================================
%  Lecture 13 -- S-GW and P-GW                            (English edition)
% =====================================================================
\documentclass[11pt,a4paper]{article}
\providecommand\ACNroot{../..}
\input{\ACNroot/templates/lecture-en}

\ACNset{number}{13}
\ACNset{title}{Serving and PDN Gateways (S-GW and P-GW)}
\ACNset{shorttitle}{S-GW and P-GW}
\ACNset{titleother}{دروازه‌های خدمت‌دهنده و PDN (S-GW و P-GW)}
\ACNset{part}{Part II -- 4G access and core (Lectures 5--13)}
\ACNset{time}{6--8 hours (3 h reading, 3--5 h exercises)}
\ACNset{level}{Intermediate}
\ACNset{prereq}{Lectures 9 and 12; IP addressing, UDP, tunnelling}
\ACNset{gen}{4}

\begin{document}
\ACNcover

% ---------------------------------------------------------------------
\section{Learning objectives}
After studying this lecture you should be able to:
\begin{objectives}
  \item \textbf{Distinguish} the functions of the serving gateway and the PDN gateway, and \textbf{explain}
        why the user plane has two anchors.
  \item \textbf{Describe} GTP-U tunnelling and \textbf{compute} its overhead and its effect on the MTU.
  \item \textbf{Trace} downlink delivery to an idle UE and the path switch after a handover.
  \item \textbf{Model} APN-AMBR policing with a token bucket and \textbf{size} gateway buffers and packet
        rates.
  \item \textbf{Explain} control/user-plane separation (CUPS) with PFCP rules, and \textbf{compare}
        home-routed and local-breakout roaming.
\end{objectives}

% ---------------------------------------------------------------------
\section{Prerequisite knowledge}
This lecture uses the bearer model of Lecture~9 (default and dedicated bearers, TFTs, QCI, AMBR) and the
MME procedures of Lecture~12 (attach, paging, service request). From networking you need IPv4/IPv6
addressing, UDP and the idea of tunnelling: carrying one packet as the payload of another.

\paragraph{Key terms} The vocabulary of this lecture with its Persian equivalents:
\begin{keyterms}
S-GW & \fa{دروازهٔ خدمت‌دهنده} & User-plane anchor for mobility between eNBs; buffers for idle UEs. \\
P-GW & \fa{دروازهٔ PDN} & Anchor towards a packet data network; allocates the UE's IP address. \\
GTP-U & \fa{پروتکل تونل‌زنی GPRS، صفحهٔ کاربر} & UDP-based tunnelling protocol for user packets \cite{ts29281}. \\
TEID & \fa{شناسهٔ نقطهٔ پایانی تونل} & 32-bit identifier of a tunnel at the receiving node. \\
PCEF & \fa{تابع اعمال سیاست و صورت‌حساب} & Enforcement point of PCC rules in the P-GW. \\
DDN & \fa{اعلان دادهٔ فروسو} & S-GW's request to the MME to page an idle UE. \\
End marker & \fa{نشانگر پایان} & GTP-U packet that marks the last packet on the old path after a switch. \\
Token bucket & \fa{سطل نشانه} & Policing model with rate $r$ and depth $b$. \\
CUPS & \fa{جداسازی صفحهٔ کنترل و کاربر} & Split of each gateway into a control and a user-plane function \cite{ts23214}. \\
PFCP & \fa{پروتکل کنترل پیش‌برد بسته} & Protocol by which the control plane installs packet rules \cite{ts29244}. \\
\end{keyterms}

% ---------------------------------------------------------------------
\section{Introduction and motivation}
Every byte that an LTE subscriber sends or receives crosses two gateways. The \emph{serving gateway}
(S-GW) faces the radio network: it stays fixed while the UE moves between eNBs and holds packets while
the UE sleeps. The \emph{PDN gateway} (P-GW) faces the Internet or the IMS: it gives the UE its IP
address and enforces the operator's policy and charging. The P-GW is the UE's point of presence in the
IP world, so it must not change during a session; the S-GW is a regional anchor and may change with
mobility \cite{ts23401}.

Gateways are where the core meets traffic volume. An MME is dimensioned in procedures per second; a
gateway is dimensioned in bits and packets per second, and its design is dominated by forwarding
performance, buffering and per-packet rule processing. This lecture follows a packet through the
gateways, quantifies the cost of tunnelling and policing, and ends with the separation of control and
user planes that leads to the 5G UPF.

% ---------------------------------------------------------------------
\section{Fundamental concepts}

\subsection{Division of labour}
\begin{table}[htbp]
\centering\small
\caption{Functions of the S-GW and the P-GW \cite{ts23401}.}\label{tab:roles}
\rowcolors{2}{acnmist}{white}
\begin{tabularx}{\linewidth}{@{}Y Y@{}}
\toprule
\rowcolor{white}S-GW & P-GW \\
\midrule
Local mobility anchor for inter-eNB handover; sends end markers & UE IP address allocation (IPv4, IPv6 prefix) \\
Downlink buffering in ECM-IDLE; triggers paging via DDN & Policy enforcement (PCEF): gating, QoS, APN-AMBR \\
Mobility anchor towards 2G/3G (S4, S12) & Downlink packet filtering with TFTs to select the bearer \\
Lawful interception & Lawful interception, deep packet inspection \\
Per-bearer accounting (e.g.\ for roaming settlement) & Online (Gy) and offline (Gz) charging \\
DSCP marking on the transport network & DSCP marking, NAT and firewall functions (operator choice) \\
\bottomrule
\end{tabularx}
\end{table}

A UE has exactly one S-GW at a time, which serves all its PDN connections; it may have several P-GWs,
one per PDN connection (for example one for the Internet APN and one for IMS). Operators often deploy
a \emph{combined S/P-GW} for non-roaming traffic, in which S5 becomes internal.

\subsection{Tunnels instead of routes}
The UE's IP address belongs to the P-GW's address pool. If routers between eNB and P-GW forwarded the
UE's packets by their IP header, every handover would require routing updates. Instead each bearer is
carried in GTP-U tunnels: S1-U between eNB and S-GW, S5/S8 between S-GW and P-GW (\cref{fig:gw}). A
handover only changes the tunnel endpoint (address and TEID) at the S-GW; the UE's address and the
P-GW never change. Each bearer has its own pair of tunnels in each direction, identified by the
receiver's TEID.

\begin{acnfigure}
  \resizebox{\linewidth}{!}{\acnfig{lecture-13/gw-functions}}
  \caption{S-GW and P-GW on the user-plane path with their main functions.}
  \label{fig:gw}
\end{acnfigure}

% ---------------------------------------------------------------------
\section{Detailed technical explanation}

\subsection{GTP-U}
GTP-U runs over UDP port 2152 \cite{ts29281}. Its mandatory header is 8 bytes: flags (version 1,
protocol type, and the E, S and PN bits that announce 4 optional bytes for an extension header, a sequence
number and an N-PDU number), the message type (255 for G-PDU, i.e.\ user data), the length of what follows
the mandatory header, and the TEID (\cref{fig:gtpu}). Signalling messages of GTP-U itself include Echo
Request/Response for path supervision, Error Indication (the receiver has no context for a TEID) and End
Marker (type 254). The inner packet (T-PDU) is not modified: the gateways see and police the UE's own IP
header.

\begin{acnfigure}
  \resizebox{\linewidth}{!}{\acnfig{lecture-13/gtpu-encap}}
  \caption{GTP-U encapsulation of a user packet and the mandatory GTP-U header.}
  \label{fig:gtpu}
\end{acnfigure}

\subsection{IP address allocation}
At attach (or PDN connection set-up) the UE requests PDN type IPv4, IPv6 or IPv4v6. For IPv4 the P-GW
returns an address in the Create Session Response, which reaches the UE in NAS; for IPv6 it allocates a
/64 prefix and announces it in Router Advertisements over the default bearer, from which the UE builds
its address. Private IPv4 addresses with carrier-grade NAT at or behind the P-GW are common, which is one
motivation for IPv6-only APNs.

\subsection{Policy enforcement}
The PCEF in the P-GW applies the PCC rules received over Gx (Lecture~10): it maps downlink packets to
bearers with DL TFTs, blocks or allows service data flows (gating), polices GBR/MBR per bearer and
APN-AMBR across the non-GBR bearers of a PDN connection (\cref{sec:tb}), counts volume and time per rating
group for charging, and marks packets with DSCP values derived from the QCI so that the transport network
can prioritise them. UE-AMBR is enforced by the eNB.

\subsection{Idle mode and handover}
When a UE goes to ECM-IDLE, the eNB releases its context and the MME asks the S-GW to release the S1-U
tunnels; the S5/S8 part stays. A downlink packet then finds no S1-U tunnel: the S-GW buffers it and
sends DDN to the MME, which pages the UE (\cref{fig:ddn}). After an X2 handover the target eNB asks the MME
for a path switch; the MME sends Modify Bearer Request to the S-GW, which switches the downlink to the
target eNB and sends one or more \emph{end marker} packets on the old path, so that the source eNB knows
when forwarding of in-flight packets is complete and the target can deliver packets in order.

\subsection{Control/user-plane separation (CUPS)}
Release~14 split each gateway into a control function (SGW-C, PGW-C) and one or more user-plane functions
(SGW-U, PGW-U), connected by the Sx interfaces \cite{ts23214}. The control function keeps S11, S5-C, Gx and
Gy; the user-plane function only forwards packets according to rules installed with PFCP \cite{ts29244}:
\begin{itemize}
  \item \textbf{PDR} (packet detection rule): which packets (TEID, UE address, SDF filter) the rule applies to;
  \item \textbf{FAR} (forwarding action rule): forward, drop, buffer or duplicate, and where to;
  \item \textbf{QER} (QoS enforcement rule): gating and rate limits;
  \item \textbf{URR} (usage reporting rule): what to count and when to report;
  \item \textbf{BAR} (buffering action rule): how to buffer for idle UEs.
\end{itemize}
User-plane functions can then be placed close to the radio network (for low latency and local traffic)
and scaled independently of the control plane. The 5G UPF and its N4 interface reuse this model
(Lecture~23).

% ---------------------------------------------------------------------
\section{Mathematical models}

\subsection{Tunnelling overhead and MTU}\label{sec:gtp}
\paragraph{Model} A user packet of $L$ bytes is carried with tunnel overhead $H$ (outer IP + UDP + GTP-U):
$H = 20+8+8 = 36$ bytes with IPv4 transport, $H = 56$ bytes with IPv6. The transport efficiency and the
largest user packet that fits a transport MTU $M$ without fragmentation are
\begin{equation}\label{eq:gtpeff}
  \eta = \frac{L}{L+H}, \qquad L_{\max} = M - H .
\end{equation}
\paragraph{Example} $L = 1400$, IPv4 transport: $\eta = 97.5\,\%$. With $M = 1500$, $L_{\max} = 1464$ bytes
(1444 with IPv6 transport), which is why operators configure the UE's MTU below 1500 or enlarge the
transport MTU.
\paragraph{Interpretation} Overhead matters for small packets: for voice it is a large fraction of the
packet (\cref{ex:volte}). Header compression (ROHC) on the radio does not help in the core, where the full
headers and the tunnel are carried.

\subsection{Downlink buffering}\label{sec:buffer}
\paragraph{Model} A source sends at rate $R$ to an idle UE; until the S1-U tunnel is restored after a
paging and service-request delay $D$, the S-GW must buffer
\begin{equation}\label{eq:buffer}
  B = R\,D .
\end{equation}
\paragraph{Example} $R = 2$\,Mbit/s, $D = 1.5$\,s: $B = 3$\,Mbit $= 375$\,kB. In practice TCP sources send
only a few packets before waiting for acknowledgements, so actual buffers are much smaller, and gateways
limit the buffer per bearer and drop beyond it.

\subsection{Policing with a token bucket}\label{sec:tb}
\paragraph{Model} A token bucket with rate $r$ and depth $b$ admits a traffic stream if, for every interval
of length $T$, the data sent $A(T)$ satisfy
\begin{equation}\label{eq:tb}
  A(T) \le b + r\,T .
\end{equation}
A source with peak rate $R > r$ starting with a full bucket can send at its peak for $T^\ast = b/(R-r)$.
\paragraph{Example} APN-AMBR $r = 10$\,Mbit/s, $b = 8$\,Mbit (1\,MB): in 3\,s at most $8 + 30 = 38$\,Mbit
(4.75\,MB) pass; a 40\,Mbit/s source can burst for $8/30 = 0.27$\,s before being limited to 10\,Mbit/s
(\cref{fig:tb}).
\paragraph{Interpretation} The rate fixes the long-term share; the depth determines how bursty traffic may
be. Too small a bucket hurts TCP (losses at every burst), too large a bucket lets short bursts exceed the
subscribed rate noticeably.

\begin{acnfigure}
  \acnfig{lecture-13/token-bucket}
  \caption{Token-bucket envelope $b + rT$ (\cref{eq:tb}) and a peak-rate source.}
  \label{fig:tb}
\end{acnfigure}

\subsection{Packet rate}\label{sec:pps}
\paragraph{Model} Forwarding cost is per packet, not per byte. A gateway carrying throughput $C$ with mean
packet size $\bar L$ handles
\begin{equation}\label{eq:pps}
  P = \frac{C}{8\,\bar L}\ \text{packets/s}.
\end{equation}
\paragraph{Example} $C = 40$\,Gbit/s, $\bar L = 800$\,bytes: $P = 6.25$\,Mpps -- each packet must be classified
against the PDRs of its session, counted and encapsulated within about 160\,ns on average per packet for a
single processing pipeline, which is why user planes use kernel-bypass packet I/O and hardware offload.

% ---------------------------------------------------------------------
\section{Network architecture and diagrams}
In a non-roaming network the S-GW and P-GW are usually co-located in regional data centres; with CUPS,
user-plane functions move towards the edge while the control functions stay central. For roaming, the
visited network provides the S-GW; the P-GW is either in the home network (\emph{home-routed}, over S8,
the usual case for data) or in the visited network (\emph{local breakout}, used e.g.\ for IMS voice in
some agreements). Home routing keeps charging and policy in the home network but lengthens the path
(\cref{exr:roaming}).

% ---------------------------------------------------------------------
\section{Step-by-step procedures}

\begin{procedurebox}{Downlink data for an idle UE}\label{proc:ddn}
\begin{steps}
  \item A downlink packet reaches the P-GW, which maps it to a bearer with the DL TFT and forwards it over
        S5 to the S-GW.
  \item The S-GW has no S1-U tunnel for the bearer: it buffers the packet and sends Downlink Data
        Notification to the MME (further packets are buffered without a new DDN).
  \item The MME pages the UE in its TAI list; the UE answers with a Service Request.
  \item The MME sets up the context in the eNB (Initial Context Setup), learns the eNB's TEID and sends
        Modify Bearer Request to the S-GW.
  \item The S-GW forwards the buffered packets to the eNB, followed by new traffic.
\end{steps}
\end{procedurebox}

\begin{acnfigure}
  \resizebox{\linewidth}{!}{\acnfig{lecture-13/ddn}}
  \caption{Network-triggered service request with S-GW buffering (\cref{proc:ddn}).}
  \label{fig:ddn}
\end{acnfigure}

\begin{procedurebox}{Path switch after an X2 handover}\label{proc:switch}
\begin{steps}
  \item During the handover the source eNB forwards in-flight downlink packets to the target eNB over X2.
  \item The target eNB sends Path Switch Request to the MME with its downlink TEIDs.
  \item The MME sends Modify Bearer Request to the S-GW, which switches the downlink tunnels to the target.
  \item The S-GW sends end-marker packets to the source eNB on the old tunnels; the source relays them to
        the target, which then delivers the forwarded packets before the new ones.
  \item The source eNB releases the UE context.
\end{steps}
\end{procedurebox}

% ---------------------------------------------------------------------
\section{Worked examples}

\begin{examplebox}{Overhead of a voice packet}\label{ex:volte}
\textbf{Problem.} An AMR-WB 12.65\,kbit/s frame produces about 33 bytes of RTP payload every 20\,ms; the UE
uses IPv6. Compute the size of the packet on S1-U (IPv4 transport) and the fraction that is headers.

\textbf{Solution.} Inner packet: $33 + 12\ (\mathrm{RTP}) + 8\ (\mathrm{UDP}) + 40\ (\mathrm{IPv6}) = 93$ bytes. On S1-U:
$93 + 36 = 129$ bytes. Tunnel overhead alone is $36/129 = 28\,\%$; all headers together are
$(129-33)/129 = 74\,\%$. At 50 packets/s the call uses $129\times8\times50 = 51.6$\,kbit/s on S1-U for 12.65\,kbit/s of
speech.
\end{examplebox}

\begin{examplebox}{Token-bucket admission}\label{ex:tb}
\textbf{Problem.} With $r = 10$\,Mbit/s and $b = 8$\,Mbit, a user downloads a 5\,MB file at 40\,Mbit/s from a
server, starting with a full bucket. How long does the download take?

\textbf{Solution.} $5$\,MB $= 40$\,Mbit. During the initial burst of $T^\ast = 0.27$\,s, $40\times0.27 = 10.7$\,Mbit pass;
the remaining $29.3$\,Mbit pass at 10\,Mbit/s in $2.93$\,s. Total $\approx 3.2$\,s, versus 1\,s unpoliced and 4\,s
with $b = 0$.
\end{examplebox}

\begin{examplebox}{Gateway packet rate}\label{ex:pps}
\textbf{Problem.} A P-GW serves $2\times10^5$ simultaneously active users with a mean of 200\,kbit/s each.
Compute throughput and packet rate for a mean packet of 800 bytes.

\textbf{Solution.} $C = 2\times10^5\times2\times10^5 = 40$\,Gbit/s; $P = 40\times10^9/6400 = 6.25$\,Mpps
(\cref{eq:pps}). Uplink packets (many small TCP acknowledgements) add to the packet rate while adding little
throughput.
\end{examplebox}

% ---------------------------------------------------------------------
\section{Technology comparison}
\begin{table}[htbp]
\centering\small
\caption{Packet gateways across generations.}\label{tab:cmp}
\rowcolors{2}{acnmist}{white}
\begin{tabularx}{\linewidth}{@{}l Y Y Y@{}}
\toprule
\rowcolor{white} & SGSN + GGSN (2G/3G) & S-GW + P-GW (4G) & UPF (+ SMF) (5G) \\
\midrule
Anchors            & GGSN (IP), SGSN (mobility) & P-GW (IP), S-GW (mobility) & PDU-session anchor UPF, intermediate UPFs \\
Tunnelling         & GTP-U (Gn, Iu-PS)          & GTP-U (S1-U, S5/S8)        & GTP-U (N3, N9) \\
Control of forwarding & Integrated              & Integrated; CUPS from Rel-14 & Always separated: PFCP over N4 \\
QoS granularity    & PDP context                & EPS bearer                 & QoS flow (QFI) \\
Local breakout / edge & Rare                    & Possible with CUPS         & Native (UL classifier, branching point) \\
\bottomrule
\end{tabularx}
\end{table}

\paragraph{One anchor or two?} Two anchors separate concerns: the S-GW hides radio mobility from the
P-GW, and the P-GW hides the PDN from the radio side. The cost is an extra hop and extra state. 5G keeps the
idea but makes it flexible: a session may have one UPF or a chain of them, chosen per session by the SMF.

% ---------------------------------------------------------------------
\section{Key takeaways}
\begin{takeawaybox}
\begin{itemize}
  \item S-GW: mobility anchor between eNBs, buffering and DDN for idle UEs, end markers.
  \item P-GW: IP anchor (address allocation), PCEF (TFT mapping, gating, AMBR, charging), SGi to the PDN.
  \item GTP-U tunnels carry bearers; overhead $H = 36$ (IPv4) or 56 (IPv6) bytes, $L_{\max} = M - H$.
  \item Buffering $B = RD$; policing $A(T) \le b + rT$; forwarding cost $P = C/(8\bar L)$ packets/s.
  \item CUPS splits gateways into control and user planes connected by PFCP rules (PDR, FAR, QER, URR, BAR)
        -- the direct ancestor of the 5G UPF.
\end{itemize}
\end{takeawaybox}

% ---------------------------------------------------------------------
\section{Exercises}

\begin{exercisebox}{1}{Which gateway?}\label{exr:roles}
Which gateway (S-GW or P-GW) changes in each case: (a) X2 handover between two eNBs of the same S-GW
area; (b) the UE moves to a region served by another S-GW; (c) the UE opens a second PDN connection to the
IMS APN?
\end{exercisebox}

\begin{exercisebox}{1}{End markers}\label{exr:endmarker}
Explain what could go wrong at the target eNB after a handover if the S-GW did not send end-marker packets.
\end{exercisebox}

\begin{exercisebox}{2}{MTU}\label{exr:mtu}
The transport network has an MTU of 1500 bytes and uses IPv6. What is the largest UE packet that avoids
fragmentation? What MTU would the transport network need to carry 1500-byte UE packets?
\end{exercisebox}

\begin{exercisebox}{2}{Buffer sizing}\label{exr:buffer}
An S-GW limits the buffer per bearer to 256\,kB. For a paging and service-request delay of 1.2\,s, up to
which source rate is no packet dropped? If $10^4$ UEs are paged simultaneously and each buffer fills, how much
memory is needed?
\end{exercisebox}

\begin{exercisebox}{3}{Home-routed roaming}\label{exr:roaming}
A subscriber roams 4000\,km from the home network and accesses a server located in the visited country.
Estimate the extra round-trip time of home routing compared with local breakout (propagation at
$2\times10^8$\,m/s, ignoring processing), and discuss why operators still prefer home routing for data.
\end{exercisebox}

\begin{exercisebox}{3}{Scaling a user plane}\label{exr:cups}
A PGW-U core forwards $1.5$\,Mpps. How many cores are needed for 100\,Gbit/s with a mean packet of 700 bytes if
cores may be loaded to at most 70\,\%? What does CUPS allow the operator to do when traffic doubles but the
number of subscribers does not?
\end{exercisebox}

\subsection*{Solutions}
\addcontentsline{toc}{subsection}{Solutions}

\begin{solutionbox}{\ref{exr:roles}}
(a) Neither: only the S1-U endpoint (eNB) changes. (b) The S-GW changes (S-GW relocation); the P-GW and IP address
stay. (c) A new P-GW (possibly) is added for the new PDN connection; the S-GW is shared.
\end{solutionbox}

\begin{solutionbox}{\ref{exr:endmarker}}
The target receives packets on two paths: forwarded over X2 from the source and directly from the S-GW. Without an
end marker it cannot know when the forwarded stream has ended; delivering new packets first would reorder the
flow (hurting TCP), and waiting by timer would add delay. The end marker closes the old path explicitly.
\end{solutionbox}

\begin{solutionbox}{\ref{exr:mtu}}
$L_{\max} = 1500 - 56 = 1444$ bytes. To carry 1500-byte UE packets the transport MTU must be at least
$1500 + 56 = 1556$ bytes (or 1540 with IPv4 transport); operators often run jumbo frames in the backhaul.
\end{solutionbox}

\begin{solutionbox}{\ref{exr:buffer}}
$R \le B/D = 256\times8\ \text{kbit}/1.2\ \text{s} = 1.71$\,Mbit/s. Worst-case memory $10^4\times256$\,kB $= 2.56$\,GB; in
practice far less because TCP sources stall after one window and most bursts are a few packets.
\end{solutionbox}

\begin{solutionbox}{\ref{exr:roaming}}
Home routing sends each packet visited $\to$ home $\to$ server (in the visited country) and back the same way: an
extra $2\times4000$\,km each direction, i.e.\ $4\times4000\ \text{km}/(2\times10^8\ \text{m/s}) = 80$\,ms of extra RTT. Operators
accept it because the home network keeps control of policy, charging, lawful interception and service access
(e.g.\ home APNs), and because roaming settlement is simpler; local breakout needs agreements on policy and
charging in the visited network.
\end{solutionbox}

\begin{solutionbox}{\ref{exr:cups}}
$P = 100\times10^9/(8\times700) = 17.9$\,Mpps; $17.9/1.5 = 11.9$ cores at full load, $11.9/0.7 = 17.0 \to 18$ cores. With
CUPS the operator adds user-plane capacity (more PGW-U instances or cores, possibly at new edge sites) without
touching the control plane, whose load depends on subscribers and procedures, not on traffic volume.
\end{solutionbox}

% ---------------------------------------------------------------------
\section{Further reading}
Gateway functions are specified in TS~23.401 \cite{ts23401}, GTPv2-C in TS~29.274 \cite{ts29274} and GTP-U in
TS~29.281 \cite{ts29281}; CUPS and PFCP in TS~23.214 \cite{ts23214} and TS~29.244 \cite{ts29244}; policy and charging
control in TS~23.203 \cite{ts23203}. Olsson et al.\ \cite{olsson2013epc} discuss gateway deployment, and
Rommer et al.\ \cite{rommer2020core} the evolution towards the 5G user plane.

\printbibliography[heading=subbibliography,title={\ACNstrRefs}]
\end{document}
